\preludeandfugue{The Room}
\contributor{Irena Peterson Pavliuk}

%Confinement — Restlessness — Absorption — Irreconcilability — Legacy — Metamorphosis

% \paragraph{TITLE:} The Room (in D minor)

\prelude

This room is lined with restlessness in my body. Every time I think
about my possible return to where I come from, something small and
defenseless shrinks inside, among the apparitions of fear, bulky and
gray. When I touch them with my finger\=they fall like towering
cardboard decorations that first seem insurmountable, but sway from
one concentrated and firm pressure\=then I know that deep inside I am
not at all that small and helpless. You remember, you never allowed
yourself to be small.

\fugue

I was born as a grown-up in this room. No, another room: the room next
door, which was even smaller. We moved into this one when dad asked
the dormitory concierge for a bigger and better room\=the whole twelve
square meters\=for a young family. I was my father's fourth
child. His first child remained a mystery: father did not recognize
him as his own and I trust my father’s words. The second one was also
not his own, but he loved the boy’s mother, who unfortunately died and
the boy was taken away by her relatives. And the third baby, my
half-sister, was taken away by my father's ex-wife after
divorce. Father said: you have a brother and a sister.

Regardless of age\=for example, I was a hundred years old at
birth\=we are a young family; dad doesn't tell my mother he loves her, dad
says I'm the prettiest baby he had ever seen.

At first, we live just with mom, without dad, in a tiny room that is
more like storage.  Nothing really fits there except a bed, a table,
and probably something else\=a closet? All the relatives who have
gathered to look at me as a baby are gently rustling in the corridor;
their voices are like wind; it is followed by a shriek: a piece of
wall plastering falls on someone tall, and someone hits their elbow
against the sink counter.

We started living with my dad only when I was seven months old. My
mother ran out of breastmilk, started feeding me from a bottle and put
me to sleep in a crib.

Hello, dear Room.

You are my compressed world, in which everything is skillfully stacked
one on top of the other, and other things forcefully shoved into the
chest of drawers, the door of which hangs on the insulating tape.

And if you don't get distracted, then go straight: upstairs to the
fifth floor. It’s 121 steps, but it depends how you count. Does the
last step count or not? I know all the ways of counting, I know how to
do it differently, some games involve multiplication and division, and
some involve running down and out into the street, as fast as you can!

My room envelops me in its creeping vines, and I do not know whether
it is a warm and familiar embrace, or a desire to grind my bones,
twist my body like a sheet in the wash. Within these walls, my body is
a flattened and damp towel hanging over the edge of the wash
basin. Warm hands are about to care for me, wash the towel in caustic
detergent\=oh I have so many spots\=pouring bleach on every spot;
later, under the caring and loving effect, the dirt will come out of
it, but the color and structure of the fibers will not be restored,
instead only a whitish-blurry circle will remain. The hands will twist
it and hang it fluttering in the wind to contemplate the world from
the height of the fifth floor.

My body flutters in the wind, attached by five clothespins to the
rope. A birch twig touches the burnt circle with its earrings, a rusty
balcony railing touches the other edge. Someone lights a cigarette
next to me and I absorb all the smoke and I absorb all the phone
conversations with swear words, simple manner of talking and ``witty''
profanities and I absorb the screech of a tram turning around at its
terminus and I hear the announcer say something about how this city is
the country’s most comfortable city to live in.
